\section{Series de potencias}
\label{sec:series_potencia}

\subsection{Sucesiones}
\label{subsec:sucesiones}

Consideremos una sucesi\'on de n\'umeros $\set{a_n}_{n=1}^{\infty}$. Decimos
que $A$ es su l\'imite si $\forall\ep>0$ existe $n_0$ tal que
$\abs{a_n-A}<\ep$ para todo $n\ge n_0$.

\begin{definicion}[Convergencia]
  \label{def:convergencia}
  Una sucesi\'on con l\'imite finito se dice \textbf{convergente}; si no
  converge es \textbf{divergente}.
\end{definicion}

\begin{definicion}[Sucesi\'on de Cauchy]
  \label{def:cauchy}
  Una sucesi\'on se llama \textbf{fundamental} o \textbf{de Cauchy} si
  $\forall\ep>0$ existe $n_0$ tal que
  \begin{equation}
    \label{eq:cauchy}
    \abs{a_n-a_m}<\ep
    \qquad\forall n,m\ge n_0 .
  \end{equation}
\end{definicion}

\begin{ejemplo}
  \label{ej:sucesion_cauchy}
  Sea $a_n=1-\dfrac{3}{n+2}$. Entonces
  \begin{equation}
    \label{eq:ejemplo_cauchy}
    \begin{aligned}
      \abs{a_n-a_m}
        &=\abs{1-\frac{3}{n+2}-\qty{1-\frac{3}{m+2}}}\\
        &=3\abs{\frac{n-m}{(n+2)(m+2)}}<\ep
    \end{aligned}
  \end{equation}
  para $n,m$ suficientemente grandes, de modo que la sucesi\'on es de Cauchy.
\end{ejemplo}

\begin{observacion}
  Las partes real e imaginaria de una sucesi\'on de Cauchy son sucesiones de
  Cauchy.
\end{observacion}

\subsection{Series}
\label{subsec:series}

Sea la serie infinita
\begin{equation}
  \label{eq:serie}
  S=a_0+a_1+a_2+\dots+a_n+\dots
\end{equation}
con $\set{a_n}$ los elementos de una sucesi\'on. Definimos las
\textbf{sumas parciales} $S_n=a_0+a_1+\dots+a_n$ y con ellas una nueva
sucesi\'on $\set{S_n}_{n=1}^{\infty}$.

Una serie infinita \textbf{converge} si $\forall\ep>0$ existe $n_0$ tal que
$\abs{S_p-S_n}<\ep$ para todos $n,p\ge n_0$; es decir, si la sucesi\'on de
sumas parciales es de Cauchy. En ese caso podemos hablar del
\textbf{residuo}:
\begin{equation}
  \label{eq:residuo}
  S=S_n+R_n,
  \qquad
  R_n=\sum_{j=n+1}^{\infty}a_j .
\end{equation}

Comparando $S$ con la serie de valores absolutos
\begin{equation}
  \label{eq:serie_absoluta}
  S_a=\abs{a_1}+\abs{a_2}+\dots+\abs{a_n}+\dots
\end{equation}
y usando la desigualdad triangular,
\begin{equation}
  \label{eq:desigualdad_triangular_serie}
  \abs{a_n+a_{n+1}+\dots+a_p}
    \le\abs{a_n}+\abs{a_{n+1}}+\dots+\abs{a_p},
\end{equation}
se ve que si $S_a$ converge entonces $S$ converge. En tal caso se dice que la
serie es \textbf{absolutamente convergente}.

\subsection{Convergencia uniforme}
\label{subsec:convergencia_uniforme}

Consideremos ahora una sucesi\'on de funciones definidas en un dominio $\Om$.
Si $\set{f_n(z)}_{n=1}^{\infty}$ converge para todo $z\in\Om$, el l\'imite
$f(z)$ es una funci\'on definida en $\Om$: dados $\ep>0$ y $z\in\Om$, existe
$n_0$ tal que $\abs{f_n(z)-f(z)}<\ep$ para $n\ge n_0$. Pero ese $n_0$ depende
en general de $z$.

\begin{definicion}[Convergencia uniforme]
  \label{def:convergencia_uniforme}
  Una sucesi\'on $\set{f_n(z)}$ converge \textbf{uniformemente} a $f(z)$ en
  $\Om$ si $\forall\ep>0$ existe un $n_0$ \emph{independiente de $z$} tal que
  \begin{equation}
    \label{eq:convergencia_uniforme}
    \abs{f_n(z)-f_m(z)}<\ep
    \qquad\forall n,m\ge n_0 .
  \end{equation}
\end{definicion}

\begin{observacion}
  La funci\'on l\'imite de una sucesi\'on uniformemente convergente de
  funciones continuas es continua.
\end{observacion}

\subsection{Caso de estudio: la serie geom\'etrica}
\label{subsec:serie_geometrica}

Sea
\begin{equation}
  \label{eq:geometrica}
  S=1+z+z^2+\dots+z^n+\dots,
  \qquad
  S_n=1+z+z^2+\dots+z^n .
\end{equation}

\begin{observacion}[El truco]
  Multiplicando por $z$ y restando,
  \begin{equation}
    \label{eq:truco_geometrica}
    \begin{aligned}
      zS_n&=z+z^2+\dots+z^{n+1}\\
      \imp S_n-zS_n&=1-z^{n+1}\\
      \imp S_n&=\frac{1-z^{n+1}}{1-z},\qquad z\neq 1 .
    \end{aligned}
  \end{equation}
\end{observacion}

Entonces
\begin{equation}
  \label{eq:limite_geometrica}
  S=\lim_{n\rightarrow\infty}S_n
   =\lim_{n\rightarrow\infty}\frac{1-z^{n+1}}{1-z},
\end{equation}
y se distinguen dos casos:
\begin{enumerate}
  \item si $\abs{z}<1$, entonces $z^{n+1}\rightarrow 0$ y
        $S=\dfrac{1}{1-z}$;
  \item si $\abs{z}>1$, entonces $\abs{z^{n+1}}\rightarrow\infty$ y la serie
        diverge.
\end{enumerate}

\begin{observacion}
  Existe pues un valor $R$ tal que hay convergencia si $\abs{z}<R$. Esto se
  generaliza a cualquier serie de potencias.
\end{observacion}

\subsection{Radio de convergencia}
\label{subsec:radio_convergencia}

\begin{teorema}[Radio de convergencia]
  \label{teo:radio_convergencia}
  Para cada serie de potencias existe un valor $R$, llamado \textbf{radio de
  convergencia}, con las siguientes propiedades:
  \begin{itemize}
    \item la serie converge absolutamente para $\abs{z}<R$, y si
          $0\le\rho<R$ la convergencia es uniforme para $\abs{z}\le\rho$;
    \item si $\abs{z}>R$ los t\'erminos de la serie no tienen cota y la serie
          diverge;
    \item en $\abs{z}<R$ la suma de la serie es una funci\'on anal\'itica; su
          derivada se obtiene diferenciando t\'ermino a t\'ermino y la serie
          de las derivadas tiene el mismo radio de convergencia.
  \end{itemize}
  Adem\'as vale la \textbf{f\'ormula de Hadamard},
  \begin{equation}
    \label{eq:hadamard}
    \frac{1}{R}=\limsup_{n\rightarrow\infty}\;\abs{a_n}^{1/n} .
  \end{equation}
\end{teorema}

\subsection{Teorema de Taylor}
\label{subsec:taylor_complejo}

\begin{teorema}[Taylor]
  \label{teo:taylor}
  Si $f(z)$ es anal\'itica en $\Om$, entonces
  \begin{equation}
    \label{eq:taylor_complejo}
    f(z)=\sum_{k=0}^{\infty}\frac{1}{k!}
      \eval{\dnv{}{z}{k}f(z)}_{z=z_0}(z-z_0)^k .
  \end{equation}
\end{teorema}

\subsection{Funciones anal\'iticas elementales}
\label{subsec:funciones_analiticas}

\subsubsection{Seno y coseno}
\label{subsubsec:seno_coseno}

Desarrollando $\sin z$ alrededor de $z=0$ con \eqref{eq:taylor_complejo},
\begin{equation}
  \label{eq:serie_seno}
  \sin z=z-\frac{1}{3!}z^3+\frac{1}{5!}z^5-\dots
    =\sum_{k=0}^{\infty}\frac{(-1)^k}{(2k+1)!}z^{2k+1},
\end{equation}
y derivando t\'ermino a t\'ermino, lo cual es l\'icito por el
Teorema~\ref{teo:radio_convergencia},
\begin{equation}
  \label{eq:serie_coseno}
  \cos z=\dv{}{z}\sin z
    =\sum_{k=0}^{\infty}\frac{(-1)^k}{(2k)!}z^{2k} .
\end{equation}

\begin{observacion}[Radio de convergencia]
  Para \eqref{eq:serie_seno} los coeficientes son
  $a_k=(-1)^k/(2k+1)!$. Como $k^k<(2k+1)!$ para $k$ grande,
  \begin{equation}
    \label{eq:cota_factorial}
    \frac{1}{(2k+1)!}<\frac{1}{k^k}
    \imp
    \qty{\frac{1}{(2k+1)!}}^{1/k}<\frac{1}{k},
  \end{equation}
  y como $\lim_{k\rightarrow\infty}1/k=0$, la f\'ormula de Hadamard
  \eqref{eq:hadamard} da $1/R=0$: la serie de $\sin z$ tiene radio de
  convergencia infinito.
\end{observacion}

\subsubsection{Relaci\'on con las funciones hiperb\'olicas}
\label{subsubsec:hiperbolicas}

Recordemos
\begin{equation}
  \label{eq:hiperbolicas}
  \sinh\al=\frac{e^{\al}-e^{-\al}}{2},
  \qquad
  \cosh\al=\frac{e^{\al}+e^{-\al}}{2},
\end{equation}
que satisfacen $\dv{}{\al}\sinh\al=\cosh\al$ y
$\dv{}{\al}\cosh\al=\sinh\al$. Desarrollando $\sinh$ alrededor de cero,
\begin{equation}
  \label{eq:serie_sinh}
  \sinh\al=\sum_{k=0}^{\infty}\frac{\al^{2k+1}}{(2k+1)!} .
\end{equation}
Por otro lado, evaluando \eqref{eq:serie_seno} en $z=i\al$ y usando
$i^{2k+1}=i(-1)^k$,
\begin{equation}
  \label{eq:seno_imaginario}
  \sin(i\al)=\sum_{k=0}^{\infty}\frac{(-1)^k}{(2k+1)!}(i\al)^{2k+1}
    =i\sum_{k=0}^{\infty}\frac{\al^{2k+1}}{(2k+1)!} ,
\end{equation}
de modo que, comparando con \eqref{eq:serie_sinh},
\begin{equation}
  \label{eq:seno_sinh}
  \sin(i\al)=i\sinh\al .
\end{equation}
An\'alogamente,
\begin{equation}
  \label{eq:coseno_cosh}
  \cos(i\al)=\sum_{k=0}^{\infty}\frac{\al^{2k}}{(2k)!}=\cosh\al .
\end{equation}

\subsubsection{La exponencial}
\label{subsubsec:exponencial}

Invirtiendo \eqref{eq:hiperbolicas} con $\al=iy$ se obtienen
\begin{equation}
  \label{eq:seno_coseno_exponencial}
  \sin y=\frac{e^{iy}-e^{-iy}}{2i},
  \qquad
  \cos y=\frac{e^{iy}+e^{-iy}}{2},
\end{equation}
y sum\'andolas, la \textbf{f\'ormula de Euler}
\begin{equation}
  \label{eq:euler}
  e^{iy}=\cos y+i\sin y .
\end{equation}

\begin{definicion}[Exponencial compleja]
  \label{def:exponencial}
  La funci\'on $e^z$ con $z\in\C$ es anal\'itica y vale su serie de Taylor,
  \begin{equation}
    \label{eq:serie_exponencial}
    e^z=\sum_{k=0}^{\infty}\frac{1}{k!}z^k .
  \end{equation}
\end{definicion}

Su radio de convergencia tambi\'en es infinito: como $k^{k/10}<k!$ para $k$
suficientemente grande,
\begin{equation}
  \label{eq:radio_exponencial}
  \qty{\frac{1}{k!}}^{1/k}<\qty{\frac{1}{k^{k/10}}}^{1/k}
    =\frac{1}{k^{1/10}}\xrightarrow[k\rightarrow\infty]{}0 .
\end{equation}

Separando \eqref{eq:serie_exponencial} para $z=iy$ en las potencias pares e
impares se recupera \eqref{eq:euler}:
\begin{equation}
  \label{eq:exponencial_euler}
  \small
  e^{iy}=\sum_{\ell=0}^{\infty}\frac{(-1)^\ell}{(2\ell)!}y^{2\ell}
    +i\sum_{\ell=0}^{\infty}\frac{(-1)^\ell}{(2\ell+1)!}y^{2\ell+1}
    =\cos y+i\sin y .
\end{equation}

En consecuencia, la representaci\'on polar \eqref{eq:polar_compleja} se
escribe de forma compacta
\begin{equation}
  \label{eq:polar_exponencial}
  z=\abs{z}\qty{\cos\vph+i\sin\vph}=\abs{z}e^{i\vph},
  \qquad \vph=\Argu(z),
\end{equation}
y para $z=x+iy$,
\begin{equation}
  \label{eq:exponencial_cartesiana}
  e^{z}=e^{x+iy}=e^{x}\qty{\cos y+i\sin y} .
\end{equation}

\begin{ejemplo}
  \label{ej:polar_exponencial}
  \begin{equation}
    \label{eq:ejemplos_polar}
    z_1=1+i=\sqrt{2}e^{i\pi/4},
    \qquad
    z_2=i=e^{i\pi/2} .
  \end{equation}
\end{ejemplo}

\subsubsection{El logaritmo}
\label{subsubsec:logaritmo}

\begin{definicion}[Logaritmo complejo]
  \label{def:logaritmo}
  \begin{equation}
    \label{eq:logaritmo}
    \ln(z)=\ln\qty{\abs{z}e^{i\vph}}=\ln\abs{z}+i\vph,
    \qquad \vph=\Argu(z).
  \end{equation}
\end{definicion}

De \eqref{eq:logaritmo} se siguen las propiedades usuales. Para el producto,
\begin{equation}
  \label{eq:log_producto}
  \begin{aligned}
    \ln(z_1z_2)
      &=\ln\qty{\abs{z_1z_2}e^{i(\vph_1+\vph_2)}}\\
      &=\ln\abs{z_1}+i\vph_1+\ln\abs{z_2}+i\vph_2\\
      &=\ln(z_1)+\ln(z_2),
  \end{aligned}
\end{equation}
y para el cociente, con $z_2\neq 0$,
\begin{equation}
  \label{eq:log_cociente}
  \begin{aligned}
    \ln\qty{\frac{z_1}{z_2}}
      &=\ln\qty{\frac{z_1\bar{z}_2}{\abs{z_2}^2}}
       =\ln(z_1)+\ln\qty{\frac{\bar{z}_2}{\abs{z_2}^2}}\\
      &=\ln(z_1)+\ln\qty{\frac{1}{\abs{z_2}}}-i\vph_2\\
      &=\ln(z_1)-\ln(z_2).
  \end{aligned}
\end{equation}

\begin{observacion}
  Como el argumento est\'a definido salvo m\'ultiplos de $2\pi$, el logaritmo
  complejo es multivaluado y $z=0$ vuelve a ser un punto rama.
\end{observacion}
